\documentclass{article}
\usepackage{amsmath}

\title{A Sample Clean Paper for Testing}
\author{Alice Test \and Bob Verify}

\begin{document}
\maketitle

\begin{abstract}
This paper presents a simple method for testing verification pipelines.
We show that our approach achieves consistent results across multiple runs.
\end{abstract}

\section{Introduction}
\label{sec:intro}

Research paper verification is important for maintaining scientific integrity.
In this work, we present a straightforward testing framework.
Our contributions are:
\begin{enumerate}
\item A clean paper structure for testing.
\item Consistent internal references.
\end{enumerate}

See Section~\ref{sec:method} for our methodology and Section~\ref{sec:results} for results.

\section{Methodology}
\label{sec:method}

We propose a three-step approach. First, we parse the document.
Second, we extract structural elements. Third, we verify consistency.

The key equation governing our approach is:
\begin{equation}
\label{eq:main}
S = \sum_{i=1}^{N} w_i \cdot l_i
\end{equation}
where $S$ is the composite score, $w_i$ are the layer weights, and $l_i$ are layer scores.

\section{Results}
\label{sec:results}

Table~\ref{tab:results} shows our main results. Our method achieves a score of 0.85
as measured by Equation~\ref{eq:main}.

\begin{table}[h]
\centering
\caption{Main results}
\label{tab:results}
\begin{tabular}{lc}
\hline
Method & Score \\
\hline
Baseline & 0.72 \\
Ours & 0.85 \\
\hline
\end{tabular}
\end{table}

\section{Conclusion}
\label{sec:conclusion}

We presented a testing framework that achieves consistent results.
Our approach scores 0.85, outperforming the baseline of 0.72.

\begin{thebibliography}{9}
\bibitem{smith2023} Smith, J. (2023). Paper verification methods. Journal of Testing, 1(1), 1-10.
\bibitem{jones2022} Jones, A. (2022). Automated research assessment. Proc. of VerifyConf, 42-50.
\end{thebibliography}

\end{document}
